% TOP_PROBLEM: {"problemId": "problem.polynomial-time-p-matrix-linear-complementarity-problem", "canonicalTitle": "Polynomial-Time P-Matrix Linear Complementarity Problem", "releaseRank": 424, "primaryDomain": "theoretical_computer_science", "primaryDomainLabel": "Theoretical computer science", "displayStatus": "open_with_solved_subcases", "releaseStatus": "open_with_solved_subcases"}
% !TeX program = lualatex
% Definition and sources only; this file is not an LLM solution attempt.
\documentclass[11pt]{article}
\usepackage[margin=25mm]{geometry}
\usepackage{fontspec}
\setmainfont{Latin Modern Roman}
\usepackage{amsmath,amssymb,mathtools}
\usepackage{unicode-math}
\setmathfont{Latin Modern Math}
\usepackage{xurl,hyperref}
\hypersetup{colorlinks=true,urlcolor=blue}
\setcounter{secnumdepth}{0}
\setlength{\parindent}{0pt}
\setlength{\parskip}{5pt}
\setlength{\emergencystretch}{3em}
\begin{document}
\section*{\texorpdfstring{Polynomial-Time P-Matrix Linear Complementarity Problem}{Polynomial-Time P-Matrix Linear Complementarity Problem}}\label{424}
\subsection{Definitions and mathematical statement}
A real square matrix $M$ is a $P$-matrix if every principal minor is strictly positive. For rational $M\in{\mathbb{Q}}^{n\times n}$ and $q\in{\mathbb{Q}}^n$, a linear complementarity solution consists of vectors $z,w\in{\mathbb{R}}^n$ with
\[
 w=Mz+q,\qquad z\ge0,\quad w\ge0,\quad z^{\mathsf T}w=0.
\]
The last condition is equivalent to $z_iw_i=0$ for all $i$. For a $P$-matrix the solution is unique. The task asks for a deterministic algorithm computing this solution in time polynomial in the binary encoding length of $(M,q)$, under the promise that $M$ is a $P$-matrix. Recognizing that promise is a separate problem. Arithmetic-operation counts with unbounded real precision do not establish a bit-complexity bound.
\par\smallskip\textit{Formulation and concept references:} \href{https://arxiv.org/abs/2008.08992}{[S1]}.

\subsection{Short English statement}
Can the unique solution of every rational $P$-matrix linear complementarity problem be found in polynomial bit time?
\subsection{Sources}
\begin{itemize}
\item[S1]A new combinatorial property of geometric unique sink orientations.
\url{https://arxiv.org/abs/2008.08992}.
\textit{Formulation source.}
\end{itemize}
\end{document}
